\section{Related Work}

\paragraph{Vision-Language Models and Thinking.}
Modern VLMs combine vision encoders with large language model backbones trained
jointly on paired image--text data \citep{liu2023visual,zhang2023multimodal}, and
visual question answering (VQA) has become a standard benchmark for evaluating their
visual grounding and language understanding capabilities \citep{antol2015vqa,goyal2016making}.
Recent \emph{thinking models}---such as DeepSeek-R1 \citep{deepseekai2025deepseekr1}
and the Qwen3 series---extend this with reinforcement-learning-trained extended
reasoning traces. Such traces improve calibration in grounded settings
\citep{yoon2025reasoningconfidence}, but can also induce systematic overconfidence
under long reasoning budgets \citep{lacombe2025dontthink}.
% We probe whether deliberate reasoning helps construct more human-like scene priors in the blind setting.

\paragraph{Language Shortcuts and Debiasing in VQA.}
Early analyses of VQA benchmarks revealed that text-only models achieved
non-trivial accuracy by exploiting question--answer regularities baked into
benchmark construction \citep{jabri2016revisiting,zhang2016yin,agrawal2018dont},
motivating debiasing and counterfactual intervention methods
\citep{cadene2019rubi,clark2019dont,teney2020learning,niu2021counterfactual}. As VLMs scaled, the problem re-emerged in a
different form: rather than benchmark design flaws, shortcuts became embedded
through large-scale pretraining, and recent audits confirm that language-only
signals remain exploitable in modern evaluations \citep{chen2024are,lee2024vlind}.
Hallucination and shortcut reliance have been widely documented in this context
\citep{ji2023survey,rohrbach2018object,xiao2021hallucination}, with VLMs
exploiting spurious linguistic cues \citep{han2024instinctive,ye2024mmspubench}
and failing compositional visual-linguistic bindings \citep{yuksekgonul2023when}.
Beyond language shortcuts, vision encoder architecture independently constrains VLM behavior: \citet{tong2024eyes} show that CLIP's representational blind spots propagate to downstream VLM failures even when the language backbone is held fixed, and \citet{mckinzie2024mm1} demonstrate that encoder choice and pre-training recipe are primary determinants of VLM behavior---both complementary failure modes to the linguistic priors we study.
These audits diagnose \emph{whether} models rely on language priors---\citet{chen2024are}
filter for vision-indispensable items, \citet{lee2024vlind} test with counterfactual
images. We take a complementary approach to understand the structure of those priors or how they compare to human expectations: removing the image
entirely and comparing model answer distributions against those of humans under the same condition.

\paragraph{Human Reference and Answer Distributions.}
AmbigQA \citep{min2020ambigqa} showed that many questions admit multiple plausible
answers and introduced distributional evaluation---scoring against the full set of
valid human responses rather than a single gold label---a design we adopt directly,
comparing model answer distributions against those of 40 humans under the same blind condition.
From the human side, dual-process accounts of judgment under uncertainty
\citep{tversky1974judgment,kahneman2011thinking} and the proactive brain framework
\citep{bar2007proactive} suggest that linguistic cues activate structured scene
predictions even without full sensory input, grounding our interpretation of human
blind responses as meaningful prior-driven behavior rather than noise.
A growing body of work evaluates AI systems against this kind of human reference
directly: \citet{lan2025mind} show that human disagreement distributions over VQA answers
reveal systematic discrepancies between model predictions and human responses
that top-1 accuracy masks, \citet{lee2023visalign} study alignment between human and model visual
judgments, and \citet{bao2026imagination} show that AI outputs systematically diverge
from the divergent and negating hypotheses humans generate freely.
Behavioral auditing work further shows that standard benchmarks miss persistent
biases that only surface under diverse input conditions
\citep{rottger2025safetyprompts,kwon2025whatllms}.
Our work extends this human-reference paradigm to visual questions under controlled
referential degradation, using a taxonomy that decomposes prior structure by
reasoning operation and entity type.

% \paragraph{Model Calibration and Overconfidence.}
% \citet{guo2017calibration} showed that modern neural networks are
% systematically overconfident---predicted probabilities do not reflect true
% correctness likelihoods. We find the same pathology in the generative VLM
% setting: blind outputs remain in a high token-confidence regime even for wrong
% committed answers, making this failure invisible to accuracy-only evaluation.
% \citet{kadavath2022language} showed LLMs can express calibrated uncertainty
% when asked to predict their own correctness; under our explicit no-image
% instruction, model confidence \emph{increases} instead---indicating committed
% hallucination rather than acknowledged uncertainty.

